\clearpage
\section{Convolutional Neural Networks (CNNs)}
\subsection*{(a) Functional Role and Architectural Justification}
The proposed architecture follows the classic hierarchical pattern of feature
extraction followed by global classification.
\begin{itemize}
 \item \textbf{Convolutional Layer ($16$ filters, $3 \times 3$)}: This layer serves as
the primary feature extractor. By utilizing a local receptive field, the layer captures
spatial correlations (edges, textures) that are lost in standard dense networks. The
use of $16$ filters allows the network to learn multiple distinct feature maps in
parallel. The "weight sharing" property is critical here, as the same $3 \times 3$
kernel is applied across the entire image, drastically reducing the parameter count
compared to a fully connected layer.
 \item \textbf{ReLU Activation}: The Rectified Linear Unit, defined as $f(x) = \max(0,
x)$, introduces the necessary non-linearity. Its primary advantage in this architecture
is the mitigation of the vanishing gradient problem, allowing for more efficient
backpropagation by maintaining a constant gradient for all positive activations.
 \item \textbf{Max-Pooling ($2 \times 2$)}: This downsampling operation reduces
the spatial resolution of the feature maps. By selecting the maximum value in each
$2 \times 2$ window, the layer provides a degree of \textit{translation invariance},
meaning the network can recognize a feature even if its position in the frame shifts
slightly. It also serves to prevent overfitting by reducing the total number of
parameters in the subsequent dense layers.
 \item \textbf{Fully Connected (FC) Layer ($100$ neurons)}: After local features are
extracted and downsampled, this layer performs "high-level" reasoning. It integrates
information from all 15,376 spatial features to form a global representation of the
input image.
 \item \textbf{Softmax Output ($10$ classes)}: The final layer maps the hidden
representations to a 10-dimensional vector. The Softmax function normalizes these
scores into probabilities ($p_i \in [0, 1]$ and $\sum p_i = 1$), which is essential for
multi-class classification.
\end{itemize}
\subsection*{(b) Detailed Spatial Dimension Analysis}
To track the transformation of the data volume, we apply the general formula for
output dimensions: $O = \frac{I - K + 2P}{S} + 1$, where $I$ is input size, $K$ is
kernel size, $P$ is padding, and $S$ is stride.
\begin{enumerate}
 \item \textbf{Input Volume}: $V_0 = 64 \times 64 \times 3$ (RGB channels).
 \item \textbf{Convolutional Layer}: Given $K=3, S=1, P=0$:
 \begin{equation*}
 O_{conv} = \frac{64 - 3 + 0}{1} + 1 = 62
 \end{equation*}
 The resulting volume is $\bm{62 \times 62 \times 16}$, where the depth reflects
the 16 filters.
 \item \textbf{ReLU Activation}: Point-wise operation; dimensions remain $\bm{62
\times 62 \times 16}$.
 \item \textbf{Max-Pooling Layer}: Given $K=2, S=2$:
 \begin{equation*}
 O_{pool} = \frac{62 - 2}{2} + 1 = 31
 \end{equation*}
 The resulting volume is $\bm{31 \times 31 \times 16}$.
 \item \textbf{Flattening}: Before the FC layer, the 3D volume is serialized into a 1D
vector:
 \begin{equation*}
 D_{flatten} = 31 \times 31 \times 16 = \bm{15,376} \text{ units.}
 \end{equation*}
\end{enumerate}
\subsection*{(c) Trainable Parameter Derivation}
The total number of parameters is the sum of weights ($\bm{W}$) and biases ($
\bm{b}$) for each layer.
\textbf{1. Convolutional Layer}:
Each of the 16 filters has a volume of $3 \times 3 \times 3$ (width $\times$ height $
\times$ input depth).
\begin{itemize}
 \item Weights: $(3 \times 3 \times 3) \times 16 = 432$
 \item Biases: $16$ (one per filter)
 \item \textbf{Layer Total}: $432 + 16 = \bm{448}$.
\end{itemize}
\textbf{2. First Fully Connected Layer}:
Every neuron in the 100-unit layer is connected to every element of the 15,376-unit
input vector.
\begin{itemize}
 \item Weights: $15,376 \times 100 = 1,537,600$
 \item Biases: $100$
 \item \textbf{Layer Total}: $1,537,600 + 100 = \bm{1,537,700}$.
\end{itemize}
\textbf{3. Output Softmax Layer}:
Maps 100 hidden neurons to 10 class scores.
\begin{itemize}
 \item Weights: $100 \times 10 = 1,000$
 \item Biases: $10$
 \item \textbf{Layer Total}: $1,000 + 10 = \bm{1,010}$.
\end{itemize}
\textbf{Total Network Parameters}:
\begin{equation*}
P_{total} = 448 + 1,537,700 + 1,010 = \bm{1,539,158}
\end{equation*}
\textit{Note: Approximately 99.9\% of the network's parameters are contained within
the first fully connected layer, highlighting the memory-intensive nature of
transitioning from convolutional feature maps to dense classification layers.}